% !TEX root = ../main.tex
\providecommand{\dimerSecFullTimePath}{\detokenize{figures/svgs/extension_secular_comparison/15_212148_dimer_fig4_T000_redfield_rwa_thermal/figures/time_all_signals_real_imag_abs_dimer_fig4_T000_redfield_rwa_thermal}}
\providecommand{\dimerSecFullFreqPath}{\detokenize{figures/svgs/extension_secular_comparison/15_212148_dimer_fig4_T000_redfield_rwa_thermal/figures/freq_all_signals_real_imag_abs_dimer_fig4_T000_redfield_rwa_thermal}}
\providecommand{\dimerNonSecTimePath}{\detokenize{figures/svgs/extension_secular_comparison/16_205651_dimer_fig4_T000_redfield_rwa_thermal/figures/time_all_signals_real_imag_abs_dimer_fig4_T000_redfield_rwa_thermal}}
\providecommand{\dimerNonSecFreqPath}{\detokenize{figures/svgs/extension_secular_comparison/16_205651_dimer_fig4_T000_redfield_rwa_thermal/figures/freq_all_signals_real_imag_abs_dimer_fig4_T000_redfield_rwa_thermal}}
\providecommand{\pentamerSecTimePath}{\detokenize{figures/svgs/extension_secular_comparison/17_230632_pentamer_ring_single/figures/time_all_signals_real_imag_abs_pentamer_ring_single}}
\providecommand{\pentamerSecFreqPath}{\detokenize{figures/svgs/extension_secular_comparison/17_230632_pentamer_ring_single/figures/freq_all_signals_real_imag_abs_pentamer_ring_single}}
\providecommand{\pentamerNonSecTimePath}{\detokenize{figures/svgs/extension_secular_comparison/17_231222_pentamer_ring_single/figures/time_all_signals_real_imag_abs_pentamer_ring_single}}
\providecommand{\pentamerNonSecFreqPath}{\detokenize{figures/svgs/extension_secular_comparison/17_231222_pentamer_ring_single/figures/freq_all_signals_real_imag_abs_pentamer_ring_single}}
\providecommand{\trimerSecTimePath}{\detokenize{figures/svgs/extension_secular_comparison/17_224610_trimer_ring_double/figures/time_all_signals_real_imag_abs_trimer_ring_double}}
\providecommand{\trimerSecFreqPath}{\detokenize{figures/svgs/extension_secular_comparison/17_224610_trimer_ring_double/figures/freq_all_signals_real_imag_abs_trimer_ring_double}}
\providecommand{\trimerNonSecTimePath}{\detokenize{figures/svgs/extension_secular_comparison/17_231158_trimer_ring_double/figures/time_all_signals_real_imag_abs_trimer_ring_double}}
\providecommand{\trimerNonSecFreqPath}{\detokenize{figures/svgs/extension_secular_comparison/17_231158_trimer_ring_double/figures/freq_all_signals_real_imag_abs_trimer_ring_double}}
\providecommand{\heptamerSingleSecTimePath}{\detokenize{figures/svgs/extension_secular_comparison/30_155009_heptamer_ring_single_secular_1e5/figures/time_all_signals_real_imag_heptamer_ring_single_secular_1e5}}
\providecommand{\heptamerSingleSecFreqPath}{\detokenize{figures/svgs/extension_secular_comparison/30_155009_heptamer_ring_single_secular_1e5/figures/freq_all_signals_real_imag_heptamer_ring_single_secular_1e5}}
\providecommand{\heptamerSingleNonSecTimePath}{\detokenize{figures/svgs/extension_secular_comparison/30_155005_heptamer_ring_single/figures/time_all_signals_real_imag_heptamer_ring_single}}
\providecommand{\heptamerSingleNonSecFreqPath}{\detokenize{figures/svgs/extension_secular_comparison/30_155005_heptamer_ring_single/figures/freq_all_signals_real_imag_heptamer_ring_single}}
\providecommand{\heptamerDoubleSecTimePath}{\detokenize{figures/svgs/extension_secular_comparison/29_223415_heptamer_ring_double_secular_1e5/figures/time_all_signals_real_imag_heptamer_ring_double_secular_1e5}}
\providecommand{\heptamerDoubleSecFreqPath}{\detokenize{figures/svgs/extension_secular_comparison/29_223415_heptamer_ring_double_secular_1e5/figures/freq_all_signals_real_imag_heptamer_ring_double_secular_1e5}}
\providecommand{\heptamerDoubleNonSecTimePath}{\detokenize{figures/svgs/extension_secular_comparison/29_223409_heptamer_ring_double/figures/time_all_signals_real_imag_heptamer_ring_double}}
\providecommand{\heptamerDoubleNonSecFreqPath}{\detokenize{figures/svgs/extension_secular_comparison/29_223409_heptamer_ring_double/figures/freq_all_signals_real_imag_heptamer_ring_double}}
\providecommand{\tridecamerSingleTimePath}{\detokenize{figures/svgs/extension_secular_comparison/19_225943_tridecamer_ring_single/figures/time_all_signals_real_imag_abs_tridecamer_ring_single}}
\providecommand{\tridecamerSingleFreqPath}{\detokenize{figures/svgs/extension_secular_comparison/19_225943_tridecamer_ring_single/figures/freq_all_signals_real_imag_abs_tridecamer_ring_single}}

\chapter{Secular and Non-Secular Comparison Panels}
\label{app:extension_secular_comparison}

This appendix records control comparisons between secular and non-secular
Redfield calculations for the already-generated dimer and ring-aggregate
examples. The purpose is conservative: before promoting any claim into the main
text, the same plot layouts are inspected side by side to check whether the
visual differences are negligible, modest, or clearly large enough to merit a
stronger discussion.

For the ring-extension examples, two secular reference choices appear. The
pentamer single-excitation and trimer double-excitation comparisons use
\(c_{\mathrm{sec}}=10^{-4}\), whereas the heptamer single- and
double-excitation comparisons use the stricter choice \(c_{\mathrm{sec}}=10^{-5}\),
matching the dimer benchmark more closely. In all cases, the corresponding
non-secular comparison runs use \(c_{\mathrm{sec}}=-1\). In the implementation,
this same dimensionless secular-cutoff parameter is passed to \texttt{QuTiP}
under the keyword \texttt{sec\_cutoff}.

The locally re-rendered comparison panels lead to the same qualitative
conclusion as before. For the present examples, secular and non-secular
propagation do not reorganize the spectra qualitatively. The dimer control
still shows the clearest visible deviations. The pentamer, trimer, and both
heptamer comparisons remain structurally very close, with differences limited
mainly to modest local intensity redistributions. The appendix placement is
therefore still the right level for the current stage of the analysis.

\section{Dimer Control Pair}
\label{app:sec:extension_secular_dimer}

The dimer pair acts as a reference because it compares the same
\(t_{\mathrm{wait}}=\SI{0}{\femto\second}\) setup already discussed in the
benchmark part of the thesis. The secular and non-secular panels agree on the
main diagonal and cross-peak structure. Only modest differences are visible in
the branch-resolved intensity balance, especially in the nonrephasing and
absorptive frequency-domain panels.

\begin{figure}[htbp]
    \centering
    \begin{minipage}[t]{0.49\linewidth}
        \centering
        \safeplaceholdergraphic[1.18\linewidth]{\dimerSecFullTimePath}
    \end{minipage}\hfill
    \begin{minipage}[t]{0.49\linewidth}
        \centering
        \safeplaceholdergraphic[1.18\linewidth]{\dimerNonSecTimePath}
    \end{minipage}
    \caption{Dimer control pair in the time domain at
    \(t_{\mathrm{wait}}=\SI{0}{\femto\second}\). Left: secular benchmark
    companion, \(c_{\mathrm{sec}}=10^{-5}\). Right: full non-secular
    Bloch--Redfield propagation, \(c_{\mathrm{sec}}=-1\). The visible
    differences are modest rather than structural.}
    \label{fig:app_extension_secular_dimer_time}
\end{figure}

\begin{figure}[htbp]
    \centering
    \begin{minipage}[t]{0.49\linewidth}
        \centering
        \safeplaceholdergraphic[1.18\linewidth]{\dimerSecFullFreqPath}
    \end{minipage}\hfill
    \begin{minipage}[t]{0.49\linewidth}
        \centering
        \safeplaceholdergraphic[1.18\linewidth]{\dimerNonSecFreqPath}
    \end{minipage}
    \caption{Dimer control pair in the frequency domain for the same runs as
    in~\autoref{fig:app_extension_secular_dimer_time}. The peak locations and
    overall panel structure are unchanged, while branch intensities vary
    slightly.}
    \label{fig:app_extension_secular_dimer_freq}
\end{figure}

\section{Ring-Aggregate Comparisons}
\label{app:sec:extension_secular_rings}

The ring examples are more exploratory than the dimer control. They ask whether
the already-generated larger-\(N\) extension panels change enough under
secularisation to justify a dedicated main-text discussion. For the pentamer
single-excitation, heptamer single-excitation, trimer double-excitation, and
heptamer double-excitation pairs, the answer remains largely negative: the
secular and non-secular spectra are visually almost the same, with only small
local intensity changes. The heptamer controls are nevertheless useful because
they extend the appendix to an intermediate system size in both the
single- and double-manifold settings while remaining computationally tractable.

\subsection{Pentamer Single-Excitation Pair}

The pentamer \(1\mathrm{ex}\) pair is the cleanest larger-\(N\) control because
both jobs use the same \(t_{\mathrm{wait}}=\SI{50}{\femto\second}\) snapshot.
The secular and non-secular panels are visually nearly indistinguishable in
both the time and frequency domains.

\begin{figure}[htbp]
    \centering
    \begin{minipage}[t]{0.49\linewidth}
        \centering
        \safeplaceholdergraphic[1.18\linewidth]{\pentamerSecTimePath}
    \end{minipage}\hfill
    \begin{minipage}[t]{0.49\linewidth}
        \centering
        \safeplaceholdergraphic[1.18\linewidth]{\pentamerNonSecTimePath}
    \end{minipage}
    \caption{Pentamer ring, single-excitation manifold, time-domain control
    pair at \(t_{\mathrm{wait}}=\SI{50}{\femto\second}\). Left:
    \(c_{\mathrm{sec}}=10^{-4}\). Right: \(c_{\mathrm{sec}}=-1\).}
    \label{fig:app_extension_secular_pentamer_time}
\end{figure}

\begin{figure}[htbp]
    \centering
    \begin{minipage}[t]{0.49\linewidth}
        \centering
        \safeplaceholdergraphic[1.18\linewidth]{\pentamerSecFreqPath}
    \end{minipage}\hfill
    \begin{minipage}[t]{0.49\linewidth}
        \centering
        \safeplaceholdergraphic[1.18\linewidth]{\pentamerNonSecFreqPath}
    \end{minipage}
    \caption{Pentamer ring, single-excitation manifold, frequency-domain
    control pair for the same runs as
    in~\autoref{fig:app_extension_secular_pentamer_time}. The visible
    differences are negligible on the scale of the full panel.}
    \label{fig:app_extension_secular_pentamer_freq}
\end{figure}

\subsection{Heptamer Single-Excitation Pair}

The heptamer \(1\mathrm{ex}\) pair extends the same comparison to a larger
single-manifold ring while using the stricter secular reference
\(c_{\mathrm{sec}}=10^{-5}\). In the locally re-rendered panels, the secular
and non-secular spectra are almost indistinguishable. The dominant peak near
\(\omega_{\mathrm{coh}} \approx \omega_{\mathrm{det}} \approx
\SI{1.65e4}{\per\centi\meter}\) and the weaker off-diagonal satellite pattern
appear in the same places in both calculations.

\begin{figure}[htbp]
    \centering
    \begin{minipage}[t]{0.49\linewidth}
        \centering
        \safeplaceholdergraphic[1.18\linewidth]{\heptamerSingleSecTimePath}
    \end{minipage}\hfill
    \begin{minipage}[t]{0.49\linewidth}
        \centering
        \safeplaceholdergraphic[1.18\linewidth]{\heptamerSingleNonSecTimePath}
    \end{minipage}
    \caption{Heptamer ring, single-excitation manifold, time-domain control
    pair at \(t_{\mathrm{wait}}=\SI{50}{\femto\second}\). Left:
    \(c_{\mathrm{sec}}=10^{-5}\). Right: \(c_{\mathrm{sec}}=-1\). In the local
    rerender, both panels show the same oscillatory structure to very good
    visual accuracy.}
    \label{fig:app_extension_secular_heptamer_single_time}
\end{figure}

\begin{figure}[htbp]
    \centering
    \begin{minipage}[t]{0.49\linewidth}
        \centering
        \safeplaceholdergraphic[1.18\linewidth]{\heptamerSingleSecFreqPath}
    \end{minipage}\hfill
    \begin{minipage}[t]{0.49\linewidth}
        \centering
        \safeplaceholdergraphic[1.18\linewidth]{\heptamerSingleNonSecFreqPath}
    \end{minipage}
    \caption{Heptamer ring, single-excitation manifold, frequency-domain
    control pair for the same runs as
    in~\autoref{fig:app_extension_secular_heptamer_single_time}. The dominant
    feature and the weaker satellite structure remain effectively unchanged
    under secularisation.}
    \label{fig:app_extension_secular_heptamer_single_freq}
\end{figure}

\subsection{Trimer Double-Excitation Pair}

The trimer \(2\mathrm{ex}\) pair gives the same message after the
double-excitation manifold is included. The denser nonrephasing structure is
already present in both calculations, and secularisation produces only small
local intensity changes.

\begin{figure}[htbp]
    \centering
    \begin{minipage}[t]{0.49\linewidth}
        \centering
        \safeplaceholdergraphic[1.18\linewidth]{\trimerSecTimePath}
    \end{minipage}\hfill
    \begin{minipage}[t]{0.49\linewidth}
        \centering
        \safeplaceholdergraphic[1.18\linewidth]{\trimerNonSecTimePath}
    \end{minipage}
    \caption{Trimer ring, double-excitation manifold, time-domain control pair
    at \(t_{\mathrm{wait}}=\SI{50}{\femto\second}\). Left:
    \(c_{\mathrm{sec}}=10^{-4}\). Right: \(c_{\mathrm{sec}}=-1\).}
    \label{fig:app_extension_secular_trimer_time}
\end{figure}

\begin{figure}[htbp]
    \centering
    \begin{minipage}[t]{0.49\linewidth}
        \centering
        \safeplaceholdergraphic[1.18\linewidth]{\trimerSecFreqPath}
    \end{minipage}\hfill
    \begin{minipage}[t]{0.49\linewidth}
        \centering
        \safeplaceholdergraphic[1.18\linewidth]{\trimerNonSecFreqPath}
    \end{minipage}
    \caption{Trimer ring, double-excitation manifold, frequency-domain control
    pair for the same runs as
    in~\autoref{fig:app_extension_secular_trimer_time}. The visible differences
    are small compared with the overall spectral structure.}
    \label{fig:app_extension_secular_trimer_freq}
\end{figure}

\subsection{Heptamer Double-Excitation Pair}

The heptamer \(2\mathrm{ex}\) pair provides an intermediate system size between
the small trimer test case and the still more demanding tridecamer limit. As
for the single-manifold heptamer control, the
comparison uses \(c_{\mathrm{sec}}=10^{-5}\) against
\(c_{\mathrm{sec}}=-1\).

In the local rerender, the double-manifold panels are again very close. The
frequency-domain peak positions and the main branch pattern remain intact, and
the time-domain oscillatory texture is preserved as well. The visible effect of
secularisation is therefore limited to modest local redistributions rather than
to a qualitative spectral change.

\begin{figure}[htbp]
    \centering
    \begin{minipage}[t]{0.49\linewidth}
        \centering
        \safeplaceholdergraphic[1.18\linewidth]{\heptamerDoubleSecTimePath}
    \end{minipage}\hfill
    \begin{minipage}[t]{0.49\linewidth}
        \centering
        \safeplaceholdergraphic[1.18\linewidth]{\heptamerDoubleNonSecTimePath}
    \end{minipage}
    \caption{Heptamer ring, double-excitation manifold, time-domain control
    pair at \(t_{\mathrm{wait}}=\SI{50}{\femto\second}\). Left: stricter
    secular comparison, \(c_{\mathrm{sec}}=10^{-5}\). Right:
    \(c_{\mathrm{sec}}=-1\). The locally re-rendered panels retain the same
    overall oscillatory texture, with only small visual differences.}
    \label{fig:app_extension_secular_heptamer_time}
\end{figure}

\begin{figure}[htbp]
    \centering
    \begin{minipage}[t]{0.49\linewidth}
        \centering
        \safeplaceholdergraphic[1.18\linewidth]{\heptamerDoubleSecFreqPath}
    \end{minipage}\hfill
    \begin{minipage}[t]{0.49\linewidth}
        \centering
        \safeplaceholdergraphic[1.18\linewidth]{\heptamerDoubleNonSecFreqPath}
    \end{minipage}
    \caption{Heptamer ring, double-excitation manifold, frequency-domain
    control pair for the same runs as
    in~\autoref{fig:app_extension_secular_heptamer_time}. The dominant peak
    positions and the overall branch layout remain stable, even though the
    branch intensities shift somewhat between the secular and non-secular
    calculations.}
    \label{fig:app_extension_secular_heptamer_freq}
\end{figure}

\subsection{Tridecamer Single-Excitation Context Panel}

At present, only the non-secular tridecamer \(1\mathrm{ex}\) snapshot is
available. It is therefore included here only as contextual background for the
planned double-manifold follow-up runs, not as an A/B claim.

\begin{figure}[htbp]
    \centering
    \begin{minipage}[t]{0.49\linewidth}
        \centering
        \safeplaceholdergraphic[1.18\linewidth]{\tridecamerSingleTimePath}
    \end{minipage}\hfill
    \begin{minipage}[t]{0.49\linewidth}
        \centering
        \safeplaceholdergraphic[1.18\linewidth]{\tridecamerSingleFreqPath}
    \end{minipage}
    \caption{Tridecamer ring, single-excitation manifold, non-secular
    reference snapshot. Left: time domain. Right: frequency domain. This panel
    serves as context until the corresponding secular and double-manifold runs
    are added.}
    \label{fig:app_extension_secular_tridecamer_context}
\end{figure}

\section{Follow-Up Runs}
\label{app:sec:extension_secular_followup}

Taken together, the heptamer single- and double-excitation comparisons extend
the appendix from the dimer benchmark to intermediate-size ring aggregates in
both the single- and double-excitation settings, while keeping the main-text
discussion unchanged. A natural further step would be a computationally
manageable larger-\(N\) double-manifold companion approaching the tridecamer
regime.
